\documentclass[12pt]{article}
\usepackage[utf8]{inputenc}
\usepackage[T1]{fontenc}
\usepackage{hyperref}
\title{Buyang, Langnian Language Resources for K-12 Curriculum Development}
\author{Compiled by Claude Sonnet 5.5 (claude-sonnet-5-5) for Andr\'as Kornai}
\date{}
\begin{document}
\maketitle

\section{Introduction}
Langnian Buyang (ISO 639-3: yln) is a Kra-Dai language of the Kra branch, spoken in Guangnan and Funing counties, Yunnan, China. Ethnologue separates it from E'ma Buyang (yzg). Wikipedia gives about 2,000 speakers for the whole Buyang cluster. Joshua Project gives figures of 300 to 400 for Langnian specifically. Search turned up almost nothing specific to Langnian.

\section{Primary Educational Resources}
No primers, textbooks or school programmes were found.

\section{Linguistic and Reference Materials}
\begin{itemize}
\item Wikipedia, ``Buyang language''. It covers the Buyang cluster generally, not Langnian alone. \url{https://en.wikipedia.org/wiki/Buyang_language}
\item Wikipedia, ``Kra languages''. \url{https://en.wikipedia.org/wiki/Kra_languages}
\item Austronesian Basic Vocabulary Database: Buyang (Ecun) word list. This is another Buyang variety, not Langnian. \url{https://abvd.eva.mpg.de/austronesian/language.php?id=693}
\end{itemize}

\section{Religious and Community Resources}
\begin{itemize}
\item Joshua Project, Langnian Buyang in China. It is a demographic profile; no scripture was confirmed. \url{https://joshuaproject.net/people_groups/19798}
\end{itemize}

\section{Digital and Audio Resources}
None found.

\section{Curriculum Development Context}
Little to nothing exists for Langnian specifically. It has very few speakers and is endangered. Any curriculum would require new documentation in cooperation with the community before teaching materials could be made.

\end{document}